% GENERATED FILE. Edit canonical lessons, then run npm run generate:books.
% canonical-source-hash: fnv1a64:ed1164ff15a08181
% canonical-lessons: IT-C536-inquilino, IT-C536-padrone-di-casa, IT-C536-arredare, IT-C536-mensola, IT-C536-comodino, IT-R536-ripasso-535-536, IT-R536-ripasso-final-536

\chapter{Una casa nuova}
\label{ch:it-una-casa-nuova}
\hlchaptermodality{536}

\begin{chapteropening}
\textbf{By the end of this chapter:} \emph{I can talk about renting and furnishing a home.}

Complete the last lesson of chapter 536: I can talk about renting and furnishing a home.
\end{chapteropening}

\section[l'inquilino]{l'inquilino — tenant}

\label{lesson:IT-C536-inquilino}



\begin{quote}
Before the new one: say the Italian for to fix, then the Italian for to paint.
\end{quote}

\subsection*{You'll want to know: l'inquilino}
\textbf{l'inquilino} — \textquotedblleft{}tenant\textquotedblright{}.

\begin{quote}
I nuovi inquilini del secondo piano sono molto gentili.
\end{quote}

\begin{quote}
\emph{The new tenants on the second floor are very nice.}
\end{quote}

A woman is \emph{l'inquilina}.

\begin{cousinweb}
\emph{inquilino} is a learned borrowing from Latin \emph{inquilīnus} \textquotedblleft{}a tenant, a lodger\textquotedblright{} (EtymDB).
\end{cousinweb}

\begin{grammarlens}[title={What you've built}]
One more word for telling, explaining and describing at B1.
\end{grammarlens}

\subsection*{Your turn}
\begin{itemize}
\raggedright
  \item \emph{Say it:} \emph{l'inquilino}
  \item \emph{Say it:} \emph{l'inquilino}, once more
  \item \emph{Say it:} a sentence of your own with \emph{l'inquilino}
  \item \emph{Recall:} say the Italian for to fix, then the Italian for to paint, then say \emph{l'inquilino} again
\end{itemize}

\begin{tcolorbox}[breakable,colback=teal!4,colframe=teal!35!black,title={Before you move on}]
What does l'inquilino mean? (\textquotedblleft{}Tenant\textquotedblright{}.) Say it once more.
\end{tcolorbox}
\section[il padrone di casa]{il padrone di casa — landlord}

\label{lesson:IT-C536-padrone-di-casa}



\begin{quote}
Before the new one: say the Italian for to paint, then the Italian for tenant.
\end{quote}

\subsection*{You'll want to know: il padrone di casa}
\textbf{il padrone di casa} — \textquotedblleft{}landlord\textquotedblright{}.

\begin{quote}
Il padrone di casa ha aumentato l'affitto.
\end{quote}

\begin{quote}
\emph{The landlord has put up the rent.}
\end{quote}

\emph{La padrona di casa} — landlady; also a host or hostess.

\begin{cousinweb}
\emph{padrone} is inherited from Latin \emph{patrōnus} \textquotedblleft{}a protector, a patron\textquotedblright{} (EtymDB), from \emph{pater} \textquotedblleft{}father\textquotedblright{}; \emph{casa} is inherited from Latin \emph{casa} \textquotedblleft{}a hut\textquotedblright{}.
\end{cousinweb}

\begin{grammarlens}[title={What you've built}]
One more word for telling, explaining and describing at B1.
\end{grammarlens}

\subsection*{Your turn}
\begin{itemize}
\raggedright
  \item \emph{Say it:} \emph{il padrone di casa}
  \item \emph{Say it:} \emph{il padrone di casa}, once more
  \item \emph{Say it:} a sentence of your own with \emph{il padrone di casa}
  \item \emph{Recall:} say the Italian for to paint, then the Italian for tenant, then say \emph{il padrone di casa} again
\end{itemize}

\begin{tcolorbox}[breakable,colback=teal!4,colframe=teal!35!black,title={Before you move on}]
What does il padrone di casa mean? (\textquotedblleft{}Landlord\textquotedblright{}.) Say it once more.
\end{tcolorbox}
\section[arredare]{arredare — to furnish}

\label{lesson:IT-C536-arredare}



\begin{quote}
Before the new one: say the Italian for tenant, then the Italian for landlord.
\end{quote}

\subsection*{You'll want to know: arredare}
\textbf{arredare} — \textquotedblleft{}to furnish\textquotedblright{}.

\begin{quote}
Abbiamo arredato la casa con mobili usati.
\end{quote}

\begin{quote}
\emph{We furnished the house with second-hand furniture.}
\end{quote}

The furniture as a whole is \emph{l'arredamento}.

\begin{cousinweb}
\emph{arredare} comes from Gothic \emph{rēdan} \textquotedblleft{}to provide\textquotedblright{}, through Medieval Latin (EtymDB).
\end{cousinweb}

\begin{grammarlens}[title={What you've built}]
One more word for telling, explaining and describing at B1.
\end{grammarlens}

\subsection*{Your turn}
\begin{itemize}
\raggedright
  \item \emph{Say it:} \emph{arredare}
  \item \emph{Say it:} \emph{arredare}, once more
  \item \emph{Say it:} a sentence of your own with \emph{arredare}
  \item \emph{Recall:} say the Italian for tenant, then the Italian for landlord, then say \emph{arredare} again
\end{itemize}

\begin{tcolorbox}[breakable,colback=teal!4,colframe=teal!35!black,title={Before you move on}]
What does arredare mean? (\textquotedblleft{}To furnish\textquotedblright{}.) Say it once more.
\end{tcolorbox}
\section[la mensola]{la mensola — shelf}

\label{lesson:IT-C536-mensola}



\begin{quote}
Before the new one: say the Italian for landlord, then the Italian for to furnish.
\end{quote}

\subsection*{You'll want to know: la mensola}
\textbf{la mensola} — \textquotedblleft{}shelf\textquotedblright{}.

\begin{quote}
Metto i libri sulla mensola.
\end{quote}

\begin{quote}
\emph{I put the books on the shelf.}
\end{quote}

A single shelf on the wall; a bookcase is \emph{la libreria}.

\begin{cousinweb}
\emph{mensola} comes from Latin \emph{mēnsula} \textquotedblleft{}a little table\textquotedblright{} (Lewis \& Short), from \emph{mēnsa} \textquotedblleft{}table\textquotedblright{}.
\end{cousinweb}

\begin{grammarlens}[title={What you've built}]
One more word for telling, explaining and describing at B1.
\end{grammarlens}

\subsection*{Your turn}
\begin{itemize}
\raggedright
  \item \emph{Say it:} \emph{la mensola}
  \item \emph{Say it:} \emph{la mensola}, once more
  \item \emph{Say it:} a sentence of your own with \emph{la mensola}
  \item \emph{Recall:} say the Italian for landlord, then the Italian for to furnish, then say \emph{la mensola} again
\end{itemize}

\begin{tcolorbox}[breakable,colback=teal!4,colframe=teal!35!black,title={Before you move on}]
What does la mensola mean? (\textquotedblleft{}Shelf\textquotedblright{}.) Say it once more.
\end{tcolorbox}
\section[il comodino]{il comodino — bedside table}

\label{lesson:IT-C536-comodino}



\begin{quote}
Before the new one: say the Italian for to furnish, then the Italian for shelf.
\end{quote}

\subsection*{You'll want to know: il comodino}
\textbf{il comodino} — \textquotedblleft{}bedside table\textquotedblright{}.

\begin{quote}
Lascio sempre il telefono sul comodino.
\end{quote}

\begin{quote}
\emph{I always leave my phone on the bedside table.}
\end{quote}

A little \emph{comò} \textquotedblleft{}chest of drawers\textquotedblright{}.

\begin{cousinweb}
\emph{comodino} is built on French \emph{commode} \textquotedblleft{}a chest of drawers\textquotedblright{} (EtymDB), from Latin \emph{commodus} \textquotedblleft{}convenient\textquotedblright{}.
\end{cousinweb}

\begin{grammarlens}[title={What you've built}]
One more word for telling, explaining and describing at B1.
\end{grammarlens}

\subsection*{Your turn}
\begin{itemize}
\raggedright
  \item \emph{Say it:} \emph{il comodino}
  \item \emph{Say it:} \emph{il comodino}, once more
  \item \emph{Say it:} a sentence of your own with \emph{il comodino}
  \item \emph{Recall:} say the Italian for to furnish, then the Italian for shelf, then say \emph{il comodino} again
\end{itemize}

\begin{tcolorbox}[breakable,colback=teal!4,colframe=teal!35!black,title={Before you move on}]
What does il comodino mean? (\textquotedblleft{}Bedside table\textquotedblright{}.) Say it once more.
\end{tcolorbox}
\section[(dialogue)]{Tempo libero e casa}

\label{lesson:IT-R536-ripasso-535-536}



\begin{quote}
Say the Italian for drill and screw.
\end{quote}

\subsection*{Your turn}
Say these aloud:

\begin{itemize}
\raggedright
  \item drill, screw, screwdriver, to fix, to paint
  \item tenant, landlord, to furnish, shelf, bedside table
\end{itemize}

\begin{tcolorbox}[breakable,colback=teal!4,colframe=teal!35!black,title={Before you move on}]
Drill? (\textbf{il trapano}.) Screw? (\textbf{la vite}.) Screwdriver? (\textbf{il cacciavite}.) To fix? (\textbf{aggiustare}.) To paint? (\textbf{imbiancare}.) Tenant? (\textbf{l'inquilino}.) Landlord? (\textbf{il padrone di casa}.) To furnish? (\textbf{arredare}.) Shelf? (\textbf{la mensola}.) Bedside table? (\textbf{il comodino}.)
\end{tcolorbox}
\section[(dialogue)]{Una casa nuova — once more}

\label{lesson:IT-R536-ripasso-final-536}



\begin{quote}
Say the Italian for tenant and landlord.
\end{quote}

\subsection*{Your turn}
\begin{itemize}
\raggedright
  \item \emph{Say it:} tenant, landlord, to furnish, shelf, bedside table
\end{itemize}

\begin{tcolorbox}[breakable,colback=teal!4,colframe=teal!35!black,title={Before you move on}]
Tenant? (\textbf{l'inquilino}.) Landlord? (\textbf{il padrone di casa}.) To furnish? (\textbf{arredare}.) Shelf? (\textbf{la mensola}.) Bedside table? (\textbf{il comodino}.)
\end{tcolorbox}
